\SourceLeaf{194}occupied with this subject. I have been at it for more than 30 years, and I do not believe
that anyone can have occupied himself more with this very $2^{\mathrm n}.$ part than I have,
although I have never published anything about it. The assumption that the sum of the 3
angles is less than $180^{0}$ leads to a geometry of its own, entirely
different from ours (Euclidean), which is wholly consistent in itself and which I have
developed to my own complete satisfaction, so that I can solve every problem in it except
that of determining a constant that cannot be found a priori. The greater
one takes this constant, the closer one comes to Euclidean geometry, and an
infinitely great value makes the two coincide. The propositions of this geometry appear in
part paradoxical and to the unpractised absurd. On closer, calm reflection, however, one finds
that they contain nothing impossible in themselves. For example, the three angles of a
triangle can become as small as one pleases, if one is allowed to take the sides large enough.
Yet the area of a triangle, however large its sides are taken, can never
exceed a certain bound or even reach it. All my efforts to find
a contradiction or inconsistency in this non-Euclidean geometry have
proved fruitless, and the only thing in it to which our understanding objects is that if
it were true, there would have to be in space a linear magnitude \textls{determinate in itself} (though unknown to us).
But it seems to me that, despite the empty verbal wisdom of the metaphysicians, we
know far too little, or nothing at all, about the true nature of space to
confuse what seems unnatural to us with what is \textls{Absolutely Impossible}. If
non-Euclidean geometry were the true one and that constant stood in some relation to
magnitudes within the range of our measurements on earth or in the heavens, then it could
be found a posteriori. I have therefore sometimes expressed in jest the wish
that Euclidean geometry were not the true one, for then we would have an absolute measure
a priori.

From a man who has shown himself to me to be a thinking mathematical mind, I do not
fear a misunderstanding of the foregoing. In any event you should regard it solely as
a private communication, which you must on no account use publicly or in a way that could
bring it\space%
